\chapter{Conclusioni}
\label{cap:conclusioni}

Questa tesi è partita da un problema concreto, quello di eseguire computazioni sui nodi di un cluster utilizzando dataset soggetti a vincoli di protezione, senza lasciare a chi sottomette le richieste la scelta di dove eseguirle, dato che gli enti esterni non hanno visibilità sull'infrastruttura e sono proprio i soggetti da cui i dati vanno protetti. Per affrontarlo si è lavorato su due fronti: sul primo è stato definito un modello che governa l'assegnazione dei task attraverso un insieme di politiche, statiche e dinamiche, dove alle condizioni di ammissibilità basate sulle proprietà dei nodi e sulla posizione dei dati, e ai criteri di costo, si affianca una politica di isolamento ispirata al modello \emph{Chinese Wall}, che tiene conto della storia delle elaborazioni su ciascun nodo; sul secondo questo modello è stato implementato in ambiente Kubernetes, sfruttando i meccanismi di estensione nativi per portare la valutazione delle politiche dalla fase di ammissione delle richieste fino alla scelta effettiva del nodo.

Una volta realizzato, il sistema è stato valutato sperimentalmente rispetto alla sua componente dinamica, l'unica il cui comportamento evolve con lo stato del cluster. Ne è emerso che il costo di questa politica non è fisso, ma nasce da un compromesso tra la disponibilità dei nodi e la reattività dello scheduling, e può quindi essere contenuto calibrando opportunamente la configurazione. Nel complesso il sistema realizza quanto si proponeva, con politiche definite centralmente e applicate in modo uniforme a ogni richiesta, e i limiti emersi riguardano le strategie adottate e il perimetro della valutazione più che l'impianto del modello.

Restano poi diversi sviluppi futuri lungo cui il lavoro potrebbe proseguire. Il più immediato nasce dalla valutazione stessa, che suggerisce di affiancare alla sanificazione periodica strategie più sofisticate, capaci di adattare la pulizia dei nodi allo stato del cluster e alla durata dei task in esecuzione anziché agire a cadenza fissa. Altre direzioni invece derivano dalle assunzioni adottate nel prototipo, che una versione più completa del sistema potrebbe gradualmente rimuovere: dare ai task una semantica propria, oltre a consentirne l'esecuzione reale al posto di un'immagine segnaposto, permetterebbe di stimarne la durata e abiliterebbe proprio le strategie di sanificazione appena accennate; la politica di isolamento, oggi riferita a ciò che è stato elaborato su un nodo, potrebbe tenere conto anche degli issuer che hanno task ancora attivi su di esso; la gestione dei dati, oggi puramente logica e limitata alla valutazione del costo del loro trasferimento, potrebbe estendersi al piano fisico, eseguendo concretamente lo spostamento o l'accesso remoto al momento della scelta del nodo; la possibilità di modificare i metadati durante l'esecuzione della pipeline aprirebbe la gestione delle modifiche concorrenti e delle decisioni che ne risultano invalidate. Infine, l'estensione a scenari multi-cluster porterebbe il sistema oltre il singolo cluster, con i problemi di coordinazione e di sincronizzazione dello stato che ne derivano. In tutti i casi si tratta di estensioni da inserire sull'infrastruttura esistente senza ridiscuterne i fondamenti, a conferma che il modello e la sua realizzazione costituiscono una base solida su cui continuare a costruire.
